\documentclass[12pt,a4paper]{article}

\usepackage[utf8]{inputenc}
\usepackage{graphicx}
\usepackage{geometry}
\geometry{a4paper, margin=2.54cm}

\title{Bank Customer Churn Prediction Using Machine Learning Techniques}
\author{Enis Çelik, Furkan Çimili, Ahmet Arif Sarı \\ Advisor: Dr.Öğr.Üy.Gökhan Göksel}
\date{April 2026}

\begin{document}

\maketitle
\begin{abstract}
This document serves as the foundation for the IEEE SIU 2026 submission and the comprehensive 35-40 page thesis report. It details the XGBoost and SHAP-based methodology applied to mitigate bank customer churn with high algorithmic fairness.
\end{abstract}

\section{Introduction}
Customer churn represents a pervasive challenge with attrition rates ranging between 10-30\%.

\section{Methodology}
We employed XGBoost, Random Forest, and Logistic Regression alongside SMOTE and ADASYN for imbalance correction. Model explanations are provided by SHAP.

\section{System Architecture}
The system architecture follows the C4 model rendered in Structurizr DSL, utilizing SQLite for data logging, FastAPI/Flask for the backend, and Streamlit for the user interface.

\section{Results}
Preliminary analysis indicates F1-scores exceeding 0.85 with minimal disparities across demographic lines.

\bibliographystyle{IEEEtran}
\bibliography{references}

\end{document}
